% SPDX-License-Identifier: Apache-2.0
\section{Limitations}
\label{sec:limitations}

This section is not a formality. The paper argues that unstated limitations are
the characteristic failure of PQC evidence, so the limitations are stated here
in the same detail as the results.

\subsection{What is not modelled}

\paragraph{Key generation.} The formalization does not model key generation. The
determinant results constrain the linear-algebraic structure of the construction,
not the distribution from which keys are drawn. A statement about the
determinant of an elementary matrix is not a statement about the security of a
key generation procedure.

\paragraph{The XOF.} The extendable output function is not modelled. Where the
construction uses XOF output, the formalization treats it as given. Any security
argument that depends on the XOF's properties is therefore outside what has been
checked.

\paragraph{T3 in its entirety.} The security, probabilistic and parser
obligations are all open, and they are open for a reason that is worth stating
plainly: their subjects are marked \texttt{TODO(define)}. There is no
formal statement to prove yet. This is not a proof debt that more tactic work
would clear; it is a modelling debt.

\subsection{What is proved, precisely}

Two tiers are closed with no remaining \texttt{sorry}: T0 (17 lemmas, purely
formal rewriting) and T1-A (16 lemmas, algebraic axioms). T1-D is nearly closed
(5 lemmas, 2 open). T2 has 3 of 10 closed. T3 has none.

The closed obligations are real results and they are machine-checked, but they
are supporting infrastructure for the construction rather than security theorems
about it. Reporting ``39 of 56 obligations closed'' without that qualification
would be the kind of number this paper is arguing against.

\subsection{The inventory is not a compliance claim}

The CBOM records which cryptographic assets are present in the source tree. It
does not establish that the inventory is complete with respect to any regulatory
definition of ``cryptographic asset'', and it does not establish compliance with
the BSI guidance or any other framework. Those are questions for a regulator and
an auditor, and a JSON file does not answer them.

The scan is source-level. A cryptographic asset introduced by a binary
dependency, generated at runtime, or reached through a code path the scanner does
not follow will not appear. The 9 components should be read as ``what a
source-level scan found'', not as ``what is in the system''.

\subsection{The formalization is of one construction}

This is a case study of one implementation of one construction. The tiering
methodology and the negative-control discipline are intended to generalise; the
specific lemma counts are not evidence about any other development.

\subsection{The negative control is not a completeness proof}

A failing negative control shows that the checker rejects a false statement. It
does not show that the checker rejects \emph{the} false statement that matters,
and it does not show that the formal statements capture the intended mathematics.
Both of those remain human judgements.

\subsection{Reproducibility has an expiry date}

The numbers in this paper were derived from the repositories at the snapshot
recorded in Appendix~\ref{app:artefacts}: \texttt{95056b4} in the product
repository and \texttt{3e04515} in the formalization repository, with
Isabelle2025-2. The scripts that derive them are shipped, so the numbers can be
recomputed; but a reader who recomputes them later will get different answers if
the development has moved. The counts are a snapshot, not a constant. This is a
real limitation and not a formality: a paper whose central claim is
``count what is actually proved'' must itself be pinned to a commit, or it
becomes the thing it criticises.

\subsection{Threats to the case-study findings}

The three failures in Section~\ref{sec:case-study} were found by one team in one
project. We have argued that they share a structural cause, and the argument is
that in each case the passing check was internal and the catching check external.
That is a claim about mechanism, not a statistical claim about frequency. We do
not know how often this happens elsewhere; we suspect often, because the
mechanism is not project-specific, but we have not measured it.
